\section{Procedimiento de verificación}

Este anexo resume el procedimiento con el que se obtienen las tablas del Anexo~A a partir de la hoja \texttt{BDD\_2026} del archivo de la ANT. Se presenta en pseudocódigo; su implementación usa solo operaciones generales de lectura, agrupación y conteo de datos. Antes de generar las tablas, el procedimiento comprueba las invariantes del paso~9: si el archivo cambiara y las cifras dejaran de coincidir, se detiene en lugar de producir resultados inconsistentes.

\begin{lstlisting}[style=codigo,caption={Procedimiento de verificación y generación de las tablas A1 a A7},label={lst:procedimiento}]
Entrada: hoja BDD_2026 del archivo de la ANT (corte al 30/09/2026)
Salida:  tablas A1 a A7

1. Cargar
   datos <- leer la hoja BDD_2026
   para cada siniestro: grave <- 1 si FALLECIDOS >= 1, si no 0
                        mes   <- mes de FECHA

2. Tabla A1 (estructura y etiqueta)
   contar filas, columnas y celdas vacías
   separar las 301 columnas de participantes (7 campos x 43) de las del evento
   participantes por siniestro: mediana y máximo
   contar graves y siniestros con 2 o más fallecidos

3. Tabla A2 (variación mensual)
   para cada mes: siniestros, graves, % de graves y % de partes de la PNE

4. Tabla A3 (calidad)
   contar ND, valor 0 y vacíos en la jerarquía de la vía
   contar el tipo del primer vehículo no identificado
   contar escrituras distintas del tipo de siniestro y horas en texto
   contar identificadores repetidos y filas duplicadas

5. Tabla A4 (contraste y fuga)
   totales del primer trimestre y de enero a mayo frente a INEC y ANT
   para cada siniestro: comparar el conteo de CONDICION_k = FALLECIDO
                        con FALLECIDOS (y LESIONADO con LESIONADOS)

6. Tabla A5 (gravedad por grupos)
   para cada variable del conjunto A y para la causa probable:
       agrupar y obtener n, graves y % de graves
       omitir grupos con menos de 30 casos (causas: menos de 100)

7. Tabla A6 (tramos temporales)
   referencia = ene-jun; ajuste = jul; verificación = ago-sep
   para cada tramo: siniestros, graves y % de graves

8. Tabla A7 (ventanas móviles)
   para m desde marzo hasta agosto:
       meses previos = ene..m; mes verificado = m+1
       para el mes verificado: siniestros, graves y % de graves

9. Verificar antes de escribir
   exigir 17.236 filas y 1.656 graves
   exigir 4.789 siniestros y 603 fallecidos en el primer trimestre
   exigir coincidencia total entre CONDICION_k y FALLECIDOS
   exigir que ninguna FECHA esté vacía
\end{lstlisting}
